% Generated from a `canvaslms quizzes export` of the practice
% exam by mkexempel.py; do not edit by hand.
\begin{otherlanguage}{swedish}
\subsection{Question 1}
\label{ep2:q1}
\emph{Learning objective 1, 1 point, choice.}

\emph{Lärandemål 1: Konstruera programmerarvänliga program med lämpliga namn och kommentarer, konsekvent språk och typografi}

\textbf{Välj ett alternativ.}

De fyra funktionerna nedan räknar alla ut samma sak och räknar rätt. Vilken är lättast att läsa och vidareutveckla för en annan programmerare, med tanke på namn, språk och typografi?

The options; the key is marked:
\begin{description}
  \item[{A}] \leavevmode\par
\begin{minted}{python}
def func(list, x):
    """Returnerar hur många av poängen i list som är minst x."""
    n = 0
    for i in list:
        if i >= x:
            n += 1
    return n
\end{minted}


    \emph{Feedback:} Inte den här: koden räknar rätt, men \texttt{func}, \texttt{x}, \texttt{n} och \texttt{i} säger inget om uppgiften, så läsaren måste gå igenom hela kroppen för att förstå den. Korta namn är inte neutrala, de flyttar arbetet till den som läser. Parametern \texttt{list} är dessutom döpt efter en typ, och inne i funktionen går det inte längre att nå Pythons inbyggda \texttt{list}.
  \item[{B \textit{(key)}}] \leavevmode\par
\begin{minted}{python}
def count_passed(scores, limit):
    """Returnerar hur många av poängen i scores som är minst limit."""
    passed = 0
    for score in scores:
        if score >= limit:
            passed += 1
    return passed
\end{minted}


    \emph{Feedback:} Rätt. Namnen säger vad koden gör, och de är skrivna på ett och samma sätt.
  \item[{C}] \leavevmode\par
\begin{minted}{python}
def raknaPassed(poängLista, limit):
    """Returnerar hur många av poängen i poängLista som är minst limit."""
    antal_ok = 0
    for score in poängLista:
        if score >= limit:
            antal_ok += 1
    return antal_ok
\end{minted}


    \emph{Feedback:} Inte den här: namnen säger en hel del, men de blandar språk (\texttt{poängLista}, \texttt{limit}, \texttt{antal\_ok}, \texttt{score}) och skrivsätt (\texttt{raknaPassed} med en stor bokstav inuti namnet, \texttt{antal\_ok} med understreck). Läsaren måste då hela tiden minnas hur just det här namnet skrivs.
  \item[{D}] \leavevmode\par
\begin{minted}{python}
def print_passed(scores, limit):
    """Returnerar hur många av poängen i scores som är minst limit."""
    total = 0
    for score in scores:
        if score >= limit:
            total += 1
    return total
\end{minted}


    \emph{Feedback:} Inte den här: det räcker inte att ett namn ser rätt ut. \texttt{print\_passed} skriver inte ut något, den returnerar ett tal, så den som läser ett anrop tror fel om vad funktionen gör. \texttt{total} låter som en summa, men är ett antal. Ett namn är ett löfte som Python inte kontrollerar.
\end{description}

The item's feedback:
\begin{description}
  \item[General feedback] \textbf{Tänk på:} alla fyra räknar rätt, så frågan handlar bara om vad namnen säger den som läser. Fråga dig för varje version: säger funktionens namn vad funktionen gör, säger variablernas namn vad de innehåller, och håller sig namnen till ett språk och ett skrivsätt?

\textbf{Prova:} terminalen hjälper inte här, eftersom alla fyra ger samma resultat. Täck över funktionens kropp och läs bara namnen. Gissa vad koden gör och kontrollera sedan gissningen mot koden.

\textbf{I kursen:} vecka 38, föreläsningen \emph{Variabler och utskrifter}, avsnitten \emph{Variabelnamn och konstanter} (namnet är till för den som läser, inte för datorn) och \emph{PEP 8} (vad beskrivande namn och en gemensam stil ger den som läser); \emph{Övning: Funktioner och variabler}, uppgiften \emph{Någon annans program} (varför funktioner gärna heter som ett verb och variabler som det de innehåller). Bedömning: rättningsmatrisen i \emph{Laboration (1) funktioner, variabler och utskrifter}, kriterierna \emph{Beskrivande namn} och \emph{Konsekvent språk/typografi}.
  \item[Feedback on a right answer] \texttt{count\_passed} säger vad funktionen gör, och \texttt{scores}, \texttt{limit}, \texttt{score} och \texttt{passed} säger vad värdena är. Allt är på engelska, med små bokstäver och understreck som PEP 8 vill, och namnen stämmer med vad koden faktiskt gör. Ett bra namn klarar alla tre sakerna samtidigt.
  \item[Feedback on a wrong answer] Gå igenom versionerna med tre frågor: Säger något namn fel om vad koden gör? Säger något namn ingenting? Blandas språk eller skrivsätt? Ett namn kan se prydligt ut och ändå brista i något av det här.
\end{description}

\subsection{Question 2}
\label{ep2:q2}
\emph{Learning objective 1, 1 point, choice.}

\emph{Lärandemål 1: Konstruera programmerarvänliga program med lämpliga namn och kommentarer, konsekvent språk och typografi}

\textbf{Välj ett alternativ.}

Du ska dokumentera funktionen nedan.

\begin{minted}{python}
def shipping_cost(weight_kg):
    ...  # Här ska den nya raden stå.
    if weight_kg <= 2:
        return 49
    return 49 + 10 * (weight_kg - 2)
\end{minted}

Den som ska anropa funktionen ska inte behöva läsa dess kropp, och beskrivningen ska synas med \texttt{help(\allowbreak{}shipping\_\allowbreak{}cost)\allowbreak{}}. Vilken rad ska stå först i funktionens kropp?

The options; the key is marked:
\begin{description}
  \item[{A}] \leavevmode\par
\begin{minted}{text}
# Returnerar fraktkostnaden i kronor för ett paket på weight_kg kilo.
\end{minted}


    \emph{Feedback:} Inte den här: en kommentar med \texttt{\#} står bara i källkoden. \texttt{help(\allowbreak{}shipping\_\allowbreak{}cost)\allowbreak{}} visar då enbart funktionshuvudet, och den som vill anropa funktionen måste läsa koden ändå. En kommentar och en docstring är inte samma sak med olika utseende.
  \item[{B}] \leavevmode\par
\begin{minted}{python}
"""Jämför weight_kg med 2 i en if-sats och räknar sedan ut svaret."""
\end{minted}


    \emph{Feedback:} Inte den här: \texttt{help} visar texten, men den återberättar hur funktionen räknar. Den som anropar funktionen ska inte behöva bry sig om hur, och en sådan text blir dessutom fel så fort du ändrar uträkningen.
  \item[{C}] \leavevmode\par
\begin{minted}{text}
"""Funktionen shipping_cost."""
\end{minted}


    \emph{Feedback:} Inte den här: \texttt{help} visar texten, men den säger inget som inte redan står i funktionens namn. En docstring är inte en ruta att bocka av.
  \item[{D \textit{(key)}}] \leavevmode\par
\begin{minted}{text}
"""Returnerar fraktkostnaden i kronor för ett paket på weight_kg kilo."""
\end{minted}


    \emph{Feedback:} Rätt. Texten syns med \texttt{help} och räcker för den som ska anropa funktionen.
\end{description}

The item's feedback:
\begin{description}
  \item[General feedback] \textbf{Tänk på:} frågan ställer två krav. Beskrivningen ska synas med \texttt{help}, och den som ska anropa funktionen ska inte behöva läsa kroppen. Pröva varje alternativ mot båda.

\textbf{Prova:} skriv in funktionen med en av raderna först och kör \texttt{help(\allowbreak{}shipping\_\allowbreak{}cost)\allowbreak{}}. Gör samma sak med de andra och jämför vad som visas. Läs sedan texten som om du bara ville anropa funktionen: vad vet du nu?

\textbf{I kursen:} vecka 38, föreläsningen \emph{Funktioner}, avsnittet \emph{Docstrings} (docstringen är till för den som ska använda funktionen utan att läsa kroppen, och den visas av \texttt{help}). Bedömning: rättningsmatrisen i \emph{Laboration (1) funktioner, variabler och utskrifter}, kriteriet \emph{Kommentarer}.
  \item[Feedback on a right answer] En docstring är en sträng först i funktionens kropp, och det är den \texttt{help} visar. Den ska säga vad anroparen får, uttryckt i parametrarna, inte hur funktionen räknar ut det. Den här säger vad anroparen får (fraktkostnaden i kronor) och vad parametern betyder (paketets vikt i kilo), så ingen behöver läsa kroppen.
  \item[Feedback on a wrong answer] Håll isär formen och innehållet. Formen prövar du med \texttt{help}. Innehållet prövar du genom att fråga: kan jag anropa funktionen rätt och förstå svaret utan att ha sett koden?
\end{description}

\subsection{Question 3}
\label{ep2:q3}
\emph{Learning objective 1, 1 point, multi-answer.}

\emph{Lärandemål 1: Konstruera programmerarvänliga program med lämpliga namn och kommentarer, konsekvent språk och typografi}

\textbf{Markera alla som stämmer.}

Du granskar en kurskamrats program.

\begin{minted}{python}
"""Skriver ut ett kvitto."""

MOMS = 0.25


def calcTotal(pris, antal):
    """Returnerar priset för antal varor med moms."""
    # Multiplicera pris med antal och lägg på moms.
    summa = pris * antal * (1 + MOMS)
    return summa


def print_receipt(price, count):
    """Skriver ut vad count varor à price kronor kostar."""
    print(f"Att betala: {calcTotal(price, count):.2f} kr")


def main():
    """Skriver ut kvittot för tre varor à 100 kronor."""
    print_receipt(100, 3)


if __name__ == "__main__":
    main()
\end{minted}

Vilka påståenden om programmet stämmer?

The options; every correct option is marked:
\begin{description}
  \item[{A \textit{(key)}}] \texttt{calcTotal} och \texttt{print\_receipt} skriver sammansatta namn på olika sätt. Enligt PEP 8 ska det vara \texttt{calc\_total}.
  \item[{B}] \texttt{MOMS} bör skrivas med små bokstäver, eftersom en funktion annars inte får läsa den.
  \item[{C \textit{(key)}}] Namnen blandar svenska och engelska, så språkvalet är inte konsekvent.
  \item[{D}] Programmet avbryts med \texttt{NameError}, eftersom \texttt{print\_receipt} anropar \texttt{calcTotal} med argument som heter något annat än parametrarna.
  \item[{E \textit{(key)}}] Kommentaren i \texttt{calcTotal} säger bara vad raden under gör, vilket koden redan säger. Den bör tas bort eller bytas mot en kommentar som svarar på varför.
\end{description}

The item's feedback:
\begin{description}
  \item[General feedback] \textbf{Allt eller inget:} frågan ger poäng bara om du har markerat alla påståenden som stämmer och inga andra. Du kan alltså ha tänkt rätt om det mesta och ändå få noll poäng. Pröva påståendena ett och ett.

\begin{itemize}
  \item \textbf{«\texttt{calcTotal} och \texttt{print\_receipt} skriver sammansatta namn på olika sätt»} Slå upp vad PEP 8 säger om funktionsnamn och jämför med de två namnen.
  \item \textbf{«\texttt{MOMS} bör skrivas med små bokstäver»} Slå upp vad PEP 8 säger om konstanter. Kör programmet för att se vad en funktion får läsa.
  \item \textbf{«Namnen blandar svenska och engelska»} Skriv upp alla namn i programmet och sortera dem efter språk.
  \item \textbf{«Programmet avbryts med \texttt{NameError}»} Kör programmet i terminalen.
  \item \textbf{«Kommentaren i \texttt{calcTotal} säger bara vad raden under gör»} Täck över kommentaren och läs raden under. Vad får du veta av kommentaren som koden inte redan säger?
\end{itemize}

\textbf{I kursen:} vecka 38, föreläsningen \emph{Variabler och utskrifter}, avsnitten \emph{Variabelnamn och konstanter} (namnet är till för den som läser, inte för datorn), \emph{Konstanter} (versalerna är ett meddelande till läsaren), \emph{Kommentarer och PEP 8} (en kommentar ska svara på varför, för det kan koden inte säga) och \emph{PEP 8} (vad beskrivande namn och en gemensam stil ger den som läser); föreläsningen \emph{Funktioner}, avsnittet \emph{Parametrar, argument och namn}. Bedömning: rättningsmatrisen i \emph{Laboration (1) funktioner, variabler och utskrifter}, kriterierna \emph{Beskrivande namn}, \emph{Konsekvent språk/typografi} och \emph{Kommentarer}.
  \item[Feedback on a right answer] \begin{itemize}
  \item \textbf{«\texttt{calcTotal} och \texttt{print\_receipt} skriver sammansatta namn på olika sätt»} Stämmer. PEP 8 säger små bokstäver och understreck för funktioner. \texttt{print\_receipt} följer det, \texttt{calcTotal} gör det inte.
  \item \textbf{«\texttt{MOMS} bör skrivas med små bokstäver»} Stämmer inte. Versaler säger läsaren att värdet är en konstant. Python bryr sig inte, och en funktion får läsa en konstant: \texttt{calcTotal} gör just det.
  \item \textbf{«Namnen blandar svenska och engelska»} Stämmer. \texttt{pris}, \texttt{antal}, \texttt{summa} och \texttt{MOMS} är svenska, medan \texttt{price}, \texttt{count}, \texttt{print\_receipt} och \texttt{calcTotal} är engelska.
  \item \textbf{«Programmet avbryts med \texttt{NameError}»} Stämmer inte. Ett argument behöver inte heta som parametern: vid anropet får \texttt{pris} det värde \texttt{price} har. Programmet skriver ut \texttt{Att betala: 375.\allowbreak{}00 kr}.
  \item \textbf{«Kommentaren i \texttt{calcTotal} säger bara vad raden under gör»} Stämmer. Kommentaren upprepar koden. En användbar kommentar svarar på det koden inte kan säga: varför.
\end{itemize}
  \item[Feedback on a wrong answer] Antingen har du markerat något som inte ska vara med, eller missat något. Skilj på det Python bryr sig om, som du avgör genom att köra programmet, och det läsaren bryr sig om, som du avgör genom att granska koden. Att ett program går att köra säger inget om hur lätt det är att läsa.
\end{description}

\subsection{Question 4}
\label{ep2:q4}
\emph{Learning objective 2, 10 points, choice.}

\emph{Lärandemål 2: Tillämpa styrstrukturer}

\textbf{Välj ett alternativ.}

En biljett kostar 40 kr för den som är under 12 år, 90 kr för den som är 12–64 år och 60 kr för den som är 65 år eller äldre. Vilken funktion returnerar rätt pris för alla åldrar?

The options; the key is marked:
\begin{description}
  \item[{A \textit{(key)}}] \leavevmode\par
\begin{minted}{python}
def ticket_price(age):
    if age >= 65:
        price = 60
    elif age >= 12:
        price = 90
    else:
        price = 40
    return price
\end{minted}


    \emph{Feedback:} Rätt. Varje ålder hamnar i rätt gren, också vid gränserna.
  \item[{B}] \leavevmode\par
\begin{minted}{python}
def ticket_price(age):
    if age < 65:
        return 90
    elif age < 12:
        return 40
    else:
        return 60
\end{minted}


    \emph{Feedback:} Inte den här: \texttt{age \textless{} 65} är sant också för barn, så \texttt{ticket\_price(5)} ger 90. Den andra grenen nås aldrig, eftersom varje tal under 12 också är under 65. Det är lätt att läsa villkoren vart och ett för sig, men i en kedja vinner det första sanna.
  \item[{C}] \leavevmode\par
\begin{minted}{python}
def ticket_price(age):
    price = 40
    if age >= 65:
        price = 60
    if age >= 12:
        price = 90
    return price
\end{minted}


    \emph{Feedback:} Inte den här: två separata \texttt{if} prövas båda, så det andra kan skriva över det första. Den som är 70 får först \texttt{price = 60} och sedan \texttt{price = 90}, så \texttt{ticket\_price(70)} ger 90.
  \item[{D}] \leavevmode\par
\begin{minted}{python}
def ticket_price(age):
    if age < 12:
        return 40
    elif age < 65 or 12:
        return 90
    else:
        return 60
\end{minted}


    \emph{Feedback:} Inte den här: \texttt{age \textless{} 65 or 12} betyder inte vad det gör på svenska. \texttt{or} prövar två delar, \texttt{age \textless{} 65} och \texttt{12}, och talet 12 räknas som sant. Hela villkoret är alltså alltid sant, grenen \texttt{else} nås aldrig och \texttt{ticket\_price(70)} ger 90.
\end{description}

The item's feedback:
\begin{description}
  \item[General feedback] \textbf{Tänk på:} följ varje version rad för rad, som datorn gör. Vilka rader körs för ett visst \texttt{age}, och vilket pris ger de?

\textbf{Prova:} anropa varje version med gränsvärdena 11, 12, 64 och 65, plus ett tal över 65, till exempel 70. Svaren ska bli 40, 90, 90, 60 och 60. Vid villkor är det gränserna som avslöjar felen.

\textbf{I kursen:} vecka 39, föreläsningen \emph{Villkor och styrstrukturer}, avsnitten \emph{Att sätta ihop villkor}, \emph{Fler än två vägar} och \emph{Tre åldersgrupper}.
  \item[Feedback on a right answer] Den som är 65 eller äldre stannar i första grenen. De som är kvar är under 65, så \texttt{age \textgreater{}= 12} skiljer ut 12–64, och resten är under 12. I en kedja med \texttt{elif} körs bara den första grenen vars villkor är sant, så varje villkor behöver bara skilja ut det som återstår. Att åldrarna prövas i omvänd ordning mot texten spelar ingen roll.
  \item[Feedback on a wrong answer] Vilket av gränsvärdena blir fel i den version du valde? Följ just det anropet rad för rad och skriv upp vilket värde \texttt{price} eller \texttt{return} får.
\end{description}

\subsection{Question 5}
\label{ep2:q5}
\emph{Learning objective 2, 10 points, choice.}

\emph{Lärandemål 2: Tillämpa styrstrukturer}

\textbf{Välj ett alternativ.}

Skriv ett litet program i terminalen som löser uppgiften, och välj sedan resultatet.

Börja med \texttt{n =\allowbreak{} 20260000000005124}. Så länge \texttt{n} är större än 10 ändras \texttt{n} enligt regeln:

\begin{itemize}
  \item om \texttt{n} är jämnt blir det nya värdet \texttt{n // 2},
  \item annars, om \texttt{n} är delbart med 3, blir det nya värdet \texttt{n // 3 + 1},
  \item annars blir det nya värdet \texttt{5 * n + 1}.
\end{itemize}

Hur många gånger ändras \texttt{n} innan det är högst 10?

The options; the key is marked:
\begin{description}
  \item[{A}] 94 gånger

    \emph{Feedback:} Inte den här: det här får du om du delar med \texttt{/} i stället för \texttt{//}. \texttt{/} ger alltid ett flyttal, och ett flyttal kan inte hålla så många siffror exakt, så \texttt{n} blir ett annat tal redan vid första delningen. Allt efter det räknas på fel tal. Skriv ut \texttt{n}: ser du en decimalpunkt eller \texttt{e+16}?
  \item[{B \textit{(key)}}] 113 gånger

    \emph{Feedback:} Rätt. Slingan, reglerna och räknaren stämmer alla.
  \item[{C}] 114 gånger

    \emph{Feedback:} Inte den här: det här får du med ett fel i gränsen eller i räknaren. Läs «större än 10» och «hur många gånger ändras» ord för ord: när ska slingan sluta, och vad ska räknas? Med \texttt{n = 91} kan du kontrollera det.
  \item[{D}] 120 gånger

    \emph{Feedback:} Inte den här: det här får du med två separata \texttt{if}, ett för jämna tal och ett för delbarhet med 3 med ett \texttt{else} efter sig. Då kan samma varv ändra \texttt{n} två gånger, men räknaren ökar bara en gång. Regeln säger «annars», och det gör att varje varv ska använda en regel.
  \item[{E}] 294 gånger

    \emph{Feedback:} Inte den här: det här får du om sista regeln blir \texttt{3 * n + 1}. Det är lätt gjort om du skriver regeln efter minnet: läs den igen och jämför tecken för tecken.
  \item[{F}] 296 gånger

    \emph{Feedback:} Inte den här: det här får du om du prövar delbarhet med 3 före jämnhet. Ett tal som är både jämnt och delbart med 3 hamnar då i fel regel, och ordningen spelar roll här. Med \texttt{n = 91} syns det redan vid 456.
\end{description}

The item's feedback:
\begin{description}
  \item[General feedback] \textbf{Tänk på:} programmet behöver en slinga som går så länge villkoret i uppgiften gäller, en regel per varv och en räknare. Läs reglerna i den ordning de står, och läs vad ordet «annars» betyder för din kod.

\textbf{Prova:} testa först på ett tal du kan följa för hand. Med \texttt{n = 91} ska följden bli 91, 456, 228, 114, 57, 20, 10, alltså 6 ändringar. Skriv ut \texttt{n} och räknaren i varje varv och jämför. Byt sedan till startvärdet i uppgiften, och se till att \texttt{n} hela tiden är ett heltal.

\textbf{Om programmet aldrig tar slut:} då är något fel i slingan, till exempel att \texttt{n} inte ändras. Stop-knappen stoppar det inte när \texttt{input} går genom en dialogruta, så tryck på Starta om miljön eller ladda om sidan. Dina filer försvinner då.

\textbf{I kursen:} vecka 39, föreläsningen \emph{Villkor och styrstrukturer}, avsnitten \emph{Fler än två vägar} och \emph{Slingor: att göra om}; vecka 40, föreläsningen \emph{Upprepningar}, avsnittet \emph{while: så länge villkoret gäller}.
  \item[Feedback on a right answer] En \texttt{while}-slinga går så länge \texttt{n} är större än 10. Varje varv använder exakt en regel, i den ordning de står, med \texttt{if}, \texttt{elif} och \texttt{else}, och räknaren ökar en gång per varv. Heltalsdivisionen \texttt{//} håller \texttt{n} som ett exakt heltal, hur stort det än är. Så här kan det se ut:

\begin{minted}{python}
def next_value(n):
    """Returnerar värdet som kommer efter n enligt regeln."""
    if n % 2 == 0:
        return n // 2
    elif n % 3 == 0:
        return n // 3 + 1
    else:
        return 5 * n + 1


def main():
    """Skriver ut hur många gånger n ändras innan det är högst 10."""
    n = 20260000000005124
    changes = 0
    while n > 10:
        n = next_value(n)
        changes += 1
    print(changes)


if __name__ == "__main__":
    main()
\end{minted}
  \item[Feedback on a wrong answer] Kör ditt program med \texttt{n = 91} och skriv ut \texttt{n} i varje varv. Det första värde som skiljer sig från följden ovan visar vilken regel som är fel. Stämmer alla värden, så titta på räknaren: stämmer antalet ändringar med följden? Prova sedan med startvärdet i uppgiften och kontrollera att \texttt{n} aldrig skrivs ut med decimalpunkt.
\end{description}

\subsection{Question 6}
\label{ep2:q6}
\emph{Learning objective 2, 10 points, choice.}

\emph{Lärandemål 2: Tillämpa styrstrukturer}

\textbf{Välj ett alternativ.}

Funktionen \texttt{has\_negative} ska returnera \texttt{True} om minst ett tal i listan \texttt{numbers} är negativt och annars \texttt{False}. Vilken version gör det för alla listor, även den tomma?

The options; the key is marked:
\begin{description}
  \item[{A}] \leavevmode\par
\begin{minted}{python}
def has_negative(numbers):
    found = False
    for number in numbers:
        if number < 0:
            found = True
        else:
            found = False
    return found
\end{minted}


    \emph{Feedback:} Inte den här: \texttt{found} skrivs över i varje varv, så bara det sista talet avgör. \texttt{has\_\allowbreak{}negative(\allowbreak{}[-{}1,\allowbreak{} 2]\allowbreak{})\allowbreak{}} ger \texttt{False}. Ett tal som inte är negativt säger inget om de tal som redan är genomgångna.
  \item[{B}] \leavevmode\par
\begin{minted}{python}
def has_negative(numbers):
    for number in numbers:
        if number < 0:
            found = True
    return found
\end{minted}


    \emph{Feedback:} Inte den här: finns inget negativt tal får \texttt{found} aldrig något värde, så sista raden avbryts med \texttt{UnboundLocalError}. Det gäller \texttt{[1, 2]} och den tomma listan. Prövar du bara listor med ett negativt tal ser allt ut att fungera.
  \item[{C}] \leavevmode\par
\begin{minted}{python}
def has_negative(numbers):
    for number in numbers:
        if number < 0:
            return True
        else:
            return False
\end{minted}


    \emph{Feedback:} Inte den här: båda grenarna returnerar, så funktionen är klar redan efter första talet. \texttt{has\_\allowbreak{}negative(\allowbreak{}[1,\allowbreak{} -{}2]\allowbreak{})\allowbreak{}} ger \texttt{False}, och den tomma listan ger \texttt{None}, eftersom slingan aldrig körs. \texttt{else} betyder här «annars, för just det här talet».
  \item[{D \textit{(key)}}] \leavevmode\par
\begin{minted}{python}
def has_negative(numbers):
    for number in numbers:
        if number < 0:
            return True
    return False
\end{minted}


    \emph{Feedback:} Rätt. Funktionen ger rätt svar för alla listor, också den tomma.
\end{description}

The item's feedback:
\begin{description}
  \item[General feedback] \textbf{Tänk på:} när vet funktionen säkert att svaret är ja? När vet den säkert att svaret är nej? Indraget avgör vilka rader som hör till slingan.

\textbf{Prova:} skriv in varje version och anropa den med \texttt{[]}, \texttt{[1, 2]}, \texttt{[-{}1, 2]} och \texttt{[1, -{}2]}. Svaren ska bli \texttt{False}, \texttt{False}, \texttt{True} och \texttt{True}. Felen syns bara på vissa listor, så pröva alla.

\textbf{I kursen:} vecka 40, föreläsningen \emph{Upprepningar}, avsnitten \emph{Vilka rader hör till slingan?} och \emph{Allt på en gång: ett gissningsspel}.
  \item[Feedback on a right answer] Svaret ja kan ges så fort ett negativt tal hittas, men svaret nej först när alla tal är genomgångna. Därför står \texttt{return True} inne i slingan och \texttt{return False} efter den. Det gäller också den tomma listan: slingan körs inget varv och svaret blir \texttt{False}.
  \item[Feedback on a wrong answer] Vilken av de fyra listorna ger fel i den version du valde? Följ just den listan varv för varv och skriv upp när funktionen är klar, och med vilket värde.
\end{description}

\subsection{Question 7}
\label{ep2:q7}
\emph{Learning objective 3, 100 points, choice.}

\emph{Lärandemål 3: Utforma och presentera användarvänliga utdata}

\textbf{Välj ett alternativ.}

Prislistan ska skrivas ut med priserna i en rak kolumn, så här:

\begin{minted}{text}
Mjölk   14.5 kr
Bröd     9.0 kr
Kaffe  109.0 kr
\end{minted}

Programmet är klart så när som på utskriften.

\begin{minted}{python}
def print_price_list(prices):
    """Skriver ut prislistan prices, en vara per rad."""
    for item, price in prices.items():
        ...  # Här ska utskriften stå.


def main():
    """Skriver ut en prislista med tre varor."""
    print_price_list({"Mjölk": 14.5, "Bröd": 9.0, "Kaffe": 109.0})


if __name__ == "__main__":
    main()
\end{minted}

Vilken rad ger exakt den utskriften?

The options; the key is marked:
\begin{description}
  \item[{A}] \leavevmode\par
\begin{minted}{python}
print(f"{item:6s}{price:.1f} kr")
\end{minted}


    \emph{Feedback:} Inte den här: namnen får en bredd, men inte priserna. Eftersom \texttt{9.0} är kortare än \texttt{109.0} hamnar decimalpunkten och \texttt{kr} på olika ställen: \texttt{Bröd  9.0 kr} och \texttt{Kaffe 109.0 kr}. Att namnkolumnen är rak räcker inte.
  \item[{B \textit{(key)}}] \leavevmode\par
\begin{minted}{python}
print(f"{item:6s}{price:6.1f} kr")
\end{minted}


    \emph{Feedback:} Rätt. Utskriften blir exakt den önskade, tecken för tecken.
  \item[{C}] \leavevmode\par
\begin{minted}{python}
print(f"{item:6s}{price:6d} kr")
\end{minted}


    \emph{Feedback:} Inte den här: \texttt{d} betyder heltal, och priserna är flyttal, så programmet avbryts med \texttt{ValueError} redan på första raden. Typen i en formatangivelse måste passa värdet.
  \item[{D}] \leavevmode\par
\begin{minted}{python}
print(f"{item}{price:6.1f} kr")
\end{minted}


    \emph{Feedback:} Inte den här: namnen är olika långa, så talet börjar på olika ställen: \texttt{Mjölk  14.5 kr} men \texttt{Bröd   9.0 kr}. En bredd räknas från där värdet börjar, så ett fält utan bredd flyttar allt som kommer efter det.
\end{description}

The item's feedback:
\begin{description}
  \item[General feedback] \textbf{Tänk på:} en formatangivelse kan ange bredd, antal decimaler och typ. Räkna tecknen i den önskade utskriften: var börjar och slutar varje kolumn, och hur många decimaler har priserna? Pythons dokumentation beskriver delarna under \emph{Format Specification Mini-Language}.

\textbf{Prova:} kör \texttt{print(\allowbreak{}f\textquotedbl{}\textbar{}\allowbreak{}\{\textquotesingle{}Bröd\textquotesingle{}:8s\}\textbar{}\allowbreak{}\{9.\allowbreak{}0:8.\allowbreak{}2f\}\textbar{}\allowbreak{}\textquotedbl{})\allowbreak{}}. Strecken visar var varje kolumn börjar och slutar. Ändra bredden och antalet decimaler tills du kan förklara vad varje del gör. Lägg sedan in raderna i programmet och jämför utskriften tecken för tecken.

\textbf{I kursen:} vecka 39, föreläsningen \emph{Inmatning och felhantering}, avsnittet \emph{Att fråga så att någon förstår} (\texttt{:.1f}); vecka 40, föreläsningen \emph{Upprepningar}, avsnittet \emph{Slingor i slingor} (\texttt{:4d}); vecka 41, föreläsningen \emph{Behållare: Uppslagslistor}, avsnittet \emph{Att skriva ut boken} (bredd på strängar).
  \item[Feedback on a right answer] \texttt{:6s} fyller ut varje namn till sex tecken, och \texttt{:6.1f} skriver priset med en decimal, högerställt, på sex tecken. Decimalpunkten hamnar då på samma ställe på varje rad, och utskriften blir exakt den önskade. Varje kolumn behöver sin egen bredd.
  \item[Feedback on a wrong answer] Kör din rad och lägg utskriften bredvid den önskade. Skiljer sig något, så är det antingen en kolumns bredd, antalet decimaler eller formatangivelsens typ. Vilket?
\end{description}

\subsection{Question 8}
\label{ep2:q8}
\emph{Learning objective 3, 100 points, multi-answer.}

\emph{Lärandemål 3: Utforma och presentera användarvänliga utdata}

\textbf{Markera alla som stämmer.}

Programmet nedan kör utan fel.

\begin{minted}{python}
"""Räknar ut ett sparande."""


def main():
    """Frågar efter tre tal och skriver ut svaret."""
    a = float(input("a:"))
    b = float(input("b"))
    c = int(input("c:"))
    print(a * (1 + b / 100) ** c)


if __name__ == "__main__":
    main()
\end{minted}

Så här ser en körning ut, med svaren 15000, 3.5 och 4:

\begin{minted}{text}
a:15000
b3.5
c:4
17212.845009374996
\end{minted}

Vilka ändringar gör körningen begripligare för den som kör programmet och aldrig ser källkoden?

The options; every correct option is marked:
\begin{description}
  \item[{A}] Byt variablerna \texttt{a}, \texttt{b} och \texttt{c} mot \texttt{deposit}, \texttt{rate} och \texttt{years}, utan att ändra någon av texterna som skrivs ut.
  \item[{B \textit{(key)}}] Låt frågorna säga vad som ska matas in och i vilken enhet, till exempel \texttt{Hur mycket sätter du in (\allowbreak{}kr)\allowbreak{}?\allowbreak{} }.
  \item[{C \textit{(key)}}] Inled med en rad som säger vad programmet räknar ut.
  \item[{D}] Flytta uträkningen till en egen funktion \texttt{final\_amount}, utan att ändra någon av texterna som skrivs ut.
  \item[{E \textit{(key)}}] Låt svaret bli en mening som säger vad talet är, med en precision som går att använda, till exempel \texttt{Efter 4 år har du 17212.\allowbreak{}85 kr.\allowbreak{}}.
\end{description}

The item's feedback:
\begin{description}
  \item[General feedback] \textbf{Allt eller inget:} frågan ger poäng bara om du har markerat alla påståenden som stämmer och inga andra. Du kan alltså ha tänkt rätt om det mesta och ändå få noll poäng. Pröva påståendena ett och ett. Skriv för varje ändring ner hur körningen skulle se ut efteråt, och jämför med körningen i frågan: är något annorlunda, och blir det begripligare för någon som aldrig sett koden?

\begin{itemize}
  \item \textbf{«Byt variablerna \texttt{a}, \texttt{b} och \texttt{c}»} Skriv ner körningen med de nya variabelnamnen.
  \item \textbf{«Låt frågorna säga vad som ska matas in och i vilken enhet»} Skriv ner körningen med sådana frågor.
  \item \textbf{«Inled med en rad som säger vad programmet räknar ut»} Skriv ner körningen med en sådan första rad.
  \item \textbf{«Flytta uträkningen till en egen funktion»} Skriv ner körningen med funktionen \texttt{final\_amount}.
  \item \textbf{«Låt svaret bli en mening som säger vad talet är»} Skriv ner körningen med en sådan sista rad.
\end{itemize}

\textbf{Prova:} gör ändringarna en i taget, kör programmet och jämför med körningen i frågan. Frågorna från \texttt{input} visas i en dialogruta.

\textbf{I kursen:} vecka 39, föreläsningen \emph{Inmatning och felhantering}, avsnittet \emph{Att fråga så att någon förstår}; vecka 38, föreläsningen \emph{Variabler och utskrifter}, avsnittet \emph{Utskrifter} (den som kör programmet ser bara utmatningen, aldrig koden). Bedömning: rättningsmatrisen i \emph{Laboration (2) inmatning, felhantering och styrstrukturer}, kriterierna \emph{Enkel inmatning} och \emph{Informativa utskrifter}.
  \item[Feedback on a right answer] \begin{itemize}
  \item \textbf{«Byt variablerna \texttt{a}, \texttt{b} och \texttt{c}»} Stämmer inte. Bättre namn är bra för programmeraren, men körningen ser exakt likadan ut.
  \item \textbf{«Låt frågorna säga vad som ska matas in och i vilken enhet»} Stämmer. Frågorna är också något programmet skriver ut. \texttt{b3.5} säger inget om vad \texttt{b} är, och med ett mellanslag sist i frågan skiljs svaret från texten.
  \item \textbf{«Inled med en rad som säger vad programmet räknar ut»} Stämmer. Den som kör programmet ser annars aldrig vad det är till för. Källkodens docstring syns inte i körningen.
  \item \textbf{«Flytta uträkningen till en egen funktion»} Stämmer inte. Det är bra uppdelning, men utskrifterna är desamma, så körningen blir inte begripligare.
  \item \textbf{«Låt svaret bli en mening som säger vad talet är»} Stämmer. \texttt{17212.\allowbreak{}845009374996} säger varken vad talet är eller i vilken enhet, och decimalerna betyder inget. \texttt{:.2f} ger \texttt{17212.85}.
\end{itemize}
  \item[Feedback on a wrong answer] Håll isär två läsare: den som läser koden och den som bara ser körningen. Frågan gäller bara den andra. Antingen har du markerat något som inte ska vara med, eller missat något.
\end{description}

\subsection{Question 9}
\label{ep2:q9}
\emph{Learning objective 3, 100 points, choice.}

\emph{Lärandemål 3: Utforma och presentera användarvänliga utdata}

\textbf{Välj ett alternativ.}

\begin{minted}{python}
class Reader:
    """En låntagare på ett bibliotek."""

    def __init__(self, name, loans):
        self.__name = name
        self.__loans = loans


def main():
    """Skriver ut en låntagare."""
    reader = Reader("Ronja", 3)
    print(reader)


if __name__ == "__main__":
    main()
\end{minted}

Programmet skriver ut något i stil med \texttt{\textless{}\_\allowbreak{}\_\allowbreak{}main\_\allowbreak{}\_\allowbreak{}.\allowbreak{}Reader object at 0x7f3a2c1e4d90\textgreater{}}. Det ska i stället skriva ut \texttt{Ronja har lånat 3 böcker.\allowbreak{}} utan att \texttt{main} ändras. Vilken metod ska läggas till i klassen?

The options; the key is marked:
\begin{description}
  \item[{A}] \leavevmode\par
\begin{minted}{python}
def to_string(self):
    return f"{self.__name} har lånat {self.__loans} böcker."
\end{minted}


    \emph{Feedback:} Inte den här: metoden fungerar om någon anropar den, men \texttt{print} letar inte efter det namnet, och \texttt{main} får inte ändras. Utskriften blir därför oförändrad.
  \item[{B}] \leavevmode\par
\begin{minted}{python}
def __str__(self):
    print(f"{self.__name} har lånat {self.__loans} böcker.")
\end{minted}


    \emph{Feedback:} Inte den här: \texttt{print} inne i metoden skriver ut texten, men metoden returnerar \texttt{None}. Texten syns alltså först, och sedan avbryts programmet med \texttt{TypeError}. Det som syns i terminalen är inte samma sak som det som lämnas tillbaka.
  \item[{C \textit{(key)}}] \leavevmode\par
\begin{minted}{python}
def __str__(self):
    return f"{self.__name} har lånat {self.__loans} böcker."
\end{minted}


    \emph{Feedback:} Rätt. \texttt{print(reader)} skriver ut precis det som ska stå, utan att \texttt{main} ändras.
  \item[{D}] \leavevmode\par
\begin{minted}{python}
def __str__(name, loans):
    return f"{name} har lånat {loans} böcker."
\end{minted}


    \emph{Feedback:} Inte den här: en metods första parameter är alltid objektet, vad den än heter. \texttt{print} skickar bara objektet, så \texttt{name} blir objektet och \texttt{loans} saknas. Programmet avbryts med \texttt{TypeError}. En metod får inte attributen som argument.
\end{description}

The item's feedback:
\begin{description}
  \item[General feedback] \textbf{Tänk på:} \texttt{main} får inte ändras, så \texttt{print(reader)} ska göra hela jobbet. Vilken metod letar \texttt{print} efter hos ett objekt, vad ska den metoden lämna tillbaka, och hur når den objektets attribut?

\textbf{Prova:} lägg in ett alternativ i klassen, kör programmet och läs vad som händer. Gör samma sak med de andra.

\textbf{I kursen:} vecka 41, föreläsningen \emph{Klasser och objekt}, avsnitten \emph{Med klasser} och \emph{Dundermetoder}; Laboration 4.
  \item[Feedback on a right answer] \texttt{print(reader)} anropar automatiskt objektets \texttt{\_\_str\_\_} och skriver ut strängen metoden returnerar: \texttt{Ronja har lånat 3 böcker.\allowbreak{}} Metoden ska alltså lämna tillbaka en sträng, inte skriva ut något. Den står i klassen, så den kommer åt \texttt{self.\_\_name} och \texttt{self.\_\_loans}, och \texttt{main} behöver inte ändras.
  \item[Feedback on a wrong answer] Kör den version du valde. Ändras utskriften? Avbryts programmet? Läs felmeddelandet noga: det säger vad Python väntade sig.
\end{description}

\subsection{Question 10}
\label{ep2:q10}
\emph{Learning objective 4, 1\,000 points, choice.}

\emph{Lärandemål 4: Överföra data mellan fil och program}

\textbf{Välj ett alternativ.}

Filen \texttt{scores.txt} innehåller tre rader och slutar med en radbrytning:

\begin{minted}{text}
12,Ronja
7,Pippi
15,Madicken
\end{minted}

Programmet nedan läser filen.

\begin{minted}{python}
def read_scores(filename):
    """Returnerar en uppslagslista från namn till poäng."""
    scores = {}
    with open(filename, "r") as file:
        for line in file:
            score, name = line.split(",")
            scores[name] = int(score)
    return scores


def main():
    """Skriver ut Ronjas poäng."""
    scores = read_scores("scores.txt")
    print(scores.get("Ronja", "saknas"))


if __name__ == "__main__":
    main()
\end{minted}

Vad händer när programmet körs?

The options; the key is marked:
\begin{description}
  \item[{A}] Programmet skriver ut \texttt{12}.

    \emph{Feedback:} Inte den här: det är vad programmeraren ville ha, men raderna ser renare ut i filen än de är när de har lästs in. Skriv ut uppslagslistan och jämför nycklarna med det \texttt{main} letar efter.
  \item[{B \textit{(key)}}] Programmet skriver ut \texttt{saknas}.

    \emph{Feedback:} Rätt. Programmet gör precis det, och det är lätt att tro något annat.
  \item[{C}] Programmet avbryts med \texttt{KeyError}.

    \emph{Feedback:} Inte den här: \texttt{KeyError} får du när du hämtar ett värde med hakparenteser. Metoden \texttt{get} fungerar annorlunda när nyckeln saknas: slå upp hur.
  \item[{D}] Programmet avbryts med \texttt{ValueError}, eftersom raden delas upp fel.

    \emph{Feedback:} Inte den här: varje rad har exakt ett kommatecken, så \texttt{split} ger två delar, och \texttt{int} får en sträng den kan omvandla. Uppdelningen avbryter alltså inte programmet.
\end{description}

The item's feedback:
\begin{description}
  \item[General feedback] \textbf{Tänk på:} följ programmet med en enda rad ur filen. Vad står exakt i \texttt{line}, i \texttt{score} och i \texttt{name}, och vilka nycklar får uppslagslistan? Vad gör sista raden i \texttt{main} med dem?

\textbf{Prova:} skapa filen med koden nedan, lägg in \texttt{read\_scores} i samma program och skriv ut hela uppslagslistan med \texttt{print(\allowbreak{}read\_\allowbreak{}scores(\allowbreak{}\textquotedbl{}scores.\allowbreak{}txt\textquotedbl{})\allowbreak{})\allowbreak{}}.

\begin{minted}{text}
with open("scores.txt", "w") as file:
    file.write("12,Ronja\n7,Pippi\n15,Madicken\n")
\end{minted}

\textbf{I kursen:} vecka 45, föreläsningen \emph{Arbeta med filer}, avsnitten \emph{Läsa och skriva} och \emph{Egna filformat}; vecka 41, föreläsningen \emph{Behållare: Uppslagslistor}.
  \item[Feedback on a right answer] \texttt{for line in file} ger raderna med radbrytningen kvar, så \texttt{split(\textquotedbl{},\textquotedbl{})} på raden med Ronja ger \texttt{\textquotedbl{}12\textquotedbl{}} och \texttt{\textquotedbl{}Ronja\textbackslash{}n\textquotedbl{}}. Nyckeln blir \texttt{\textquotedbl{}Ronja\textbackslash{}n\textquotedbl{}}, inte \texttt{\textquotedbl{}Ronja\textquotedbl{}}, och \texttt{get} returnerar då reservvärdet \texttt{\textquotedbl{}saknas\textquotedbl{}}. Lösningen är att ta bort radbrytningen först: \texttt{line.\allowbreak{}strip(\allowbreak{})\allowbreak{}.\allowbreak{}split(\allowbreak{}\textquotedbl{},\allowbreak{}\textquotedbl{})\allowbreak{}}. Programmet kraschar inte, och det gör felet lätt att missa. Läs mer i föreläsningen \emph{Behållare: Uppslagslistor}, avsnittet \emph{När nyckeln saknas}.
  \item[Feedback on a wrong answer] Skriv ut uppslagslistan och läs varje nyckel tecken för tecken. Jämför sedan med strängen som \texttt{main} slår upp, och läs vad dokumentationen säger om metoden \texttt{get}.
\end{description}

\subsection{Question 11}
\label{ep2:q11}
\emph{Learning objective 4, 1\,000 points, choice.}

\emph{Lärandemål 4: Överföra data mellan fil och program}

\textbf{Välj ett alternativ.}

Ett program håller reda på ett lager i en uppslagslista, till exempel \texttt{\{\textquotedbl{}mjölk\textquotedbl{}: 12,\allowbreak{} \textquotedbl{}bröd\textquotedbl{}: 7\}}. Funktionen \texttt{save\_stock} anropas varje gång lagret ändrats. Filen ska efter varje anrop innehålla exakt det som finns i uppslagslistan just då, en vara per rad, på formatet \texttt{mjölk,12}, så att programmet kan läsa tillbaka den nästa gång det startas. Vilken version gör det?

The options; the key is marked:
\begin{description}
  \item[{A \textit{(key)}}] \leavevmode\par
\begin{minted}{python}
def save_stock(stock, filename):
    with open(filename, "w") as file:
        for item, count in stock.items():
            file.write(f"{item},{count}\n")
\end{minted}


    \emph{Feedback:} Rätt. Filen innehåller efter varje anrop exakt det som finns just då, en vara per rad.
  \item[{B}] \leavevmode\par
\begin{minted}{python}
def save_stock(stock, filename):
    with open(filename, "a") as file:
        for item, count in stock.items():
            file.write(f"{item},{count}\n")
\end{minted}


    \emph{Feedback:} Inte den här: läget \texttt{\textquotedbl{}a\textquotedbl{}} lägger till sist i filen, så de gamla raderna står kvar. Efter två anrop kan samma vara finnas två gånger, till exempel \texttt{mjölk,12} och \texttt{mjölk,11}. Det läget passar för en logg, inte för en fil som ska visa läget just nu.
  \item[{C}] \leavevmode\par
\begin{minted}{python}
def save_stock(stock, filename):
    with open(filename, "w") as file:
        for item, count in stock.items():
            file.write(f"{item},{count}")
\end{minted}


    \emph{Feedback:} Inte den här: \texttt{write} lägger inte till någon radbrytning, som \texttt{print} gör, så allt hamnar på en rad: \texttt{mjölk,12bröd,7}. Sedan går det inte att se var en vara slutar och nästa börjar.
  \item[{D}] \leavevmode\par
\begin{minted}{python}
def save_stock(stock, filename):
    with open(filename, "w") as file:
        for item, count in stock.items():
            file.write(item + "," + count + "\n")
\end{minted}


    \emph{Feedback:} Inte den här: \texttt{count} är ett \texttt{int}, och \texttt{+} mellan en sträng och ett heltal ger \texttt{TypeError}. Eftersom \texttt{\textquotedbl{}w\textquotedbl{}} redan tömt filen är det gamla lagret dessutom borta när felet kommer.
\end{description}

The item's feedback:
\begin{description}
  \item[General feedback] \textbf{Tänk på:} kraven är tre. Filen ska efter varje anrop innehålla exakt det som finns just då, varje vara ska stå på en egen rad, och programmet ska kunna läsa tillbaka den. Pröva varje version mot alla tre, också vid det andra anropet.

\textbf{Prova:} lägg in en version av \texttt{save\_stock} i ett program tillsammans med koden nedan och kör det. Efteråt ska filen innehålla exakt en rad, \texttt{mjölk,11}.

\begin{minted}{python}
save_stock({"mjölk": 12, "bröd": 7}, "stock.txt")
save_stock({"mjölk": 11}, "stock.txt")
with open("stock.txt", "r") as file:
    print(file.read())
\end{minted}

\textbf{I kursen:} vecka 45, föreläsningen \emph{Arbeta med filer}, avsnittet \emph{Läsa och skriva}; vecka 38, föreläsningen \emph{Variabler och utskrifter}, avsnittet \emph{Att konvertera mellan typer}.
  \item[Feedback on a right answer] Läget \texttt{\textquotedbl{}w\textquotedbl{}} tömmer filen när den öppnas, så inget gammalt står kvar. Varje vara får en egen rad tack vare \texttt{\textbackslash{}n}, som måste stå med eftersom \texttt{write}, till skillnad från \texttt{print}, inte lägger till någon radbrytning. F-strängen gör om \texttt{count} till text. Filen får raderna \texttt{mjölk,12} och \texttt{bröd,7}, som går att läsa tillbaka rad för rad.
  \item[Feedback on a wrong answer] Kör din version med två anrop och läs filen. Står något kvar från första anropet? Står varje vara på en egen rad? Avbröts programmet?
\end{description}

\subsection{Question 12}
\label{ep2:q12}
\emph{Learning objective 4, 1\,000 points, multi-answer.}

\emph{Lärandemål 4: Överföra data mellan fil och program}

\textbf{Markera alla som stämmer.}

Filen \texttt{results.csv} innehåller:

\begin{minted}{text}
Birk,3
"Rövardotter, Ronja",5
\end{minted}

Här är två sätt att läsa den. Först med \texttt{split}:

\begin{minted}{python}
with open("results.csv", "r") as file:
    for line in file:
        name, score = line.strip().split(",")
\end{minted}

och sedan med modulen \texttt{csv}:

\begin{minted}{python}
import csv

with open("results.csv", "r", newline="") as file:
    for name, score in csv.reader(file):
        ...
\end{minted}

Vilka påståenden stämmer?

The options; every correct option is marked:
\begin{description}
  \item[{A}] Filen behöver inte öppnas om modulen \texttt{csv} används: \texttt{csv.\allowbreak{}reader(\allowbreak{}\textquotedbl{}results.\allowbreak{}csv\textquotedbl{})\allowbreak{}} läser filen med det namnet.
  \item[{B \textit{(key)}}] Läsningen med \texttt{split} avbryts med \texttt{ValueError} på raden med Ronja, eftersom raden delas i tre delar.
  \item[{C}] Med modulen \texttt{csv} är \texttt{score} ett heltal, eftersom modulen omvandlar fält som ser ut som tal.
  \item[{D \textit{(key)}}] Med modulen \texttt{csv} får \texttt{name} på raden med Ronja värdet \texttt{Rövardotter,\allowbreak{} Ronja}, utan citattecken.
  \item[{E}] Efter läsningen med modulen \texttt{csv} måste \texttt{file.close()} anropas, eftersom \texttt{with} inte stänger en fil som \texttt{csv.reader} har läst.
\end{description}

The item's feedback:
\begin{description}
  \item[General feedback] \textbf{Allt eller inget:} frågan ger poäng bara om du har markerat alla påståenden som stämmer och inga andra. Du kan alltså ha tänkt rätt om det mesta och ändå få noll poäng. Pröva påståendena ett och ett. Alla fem går att pröva i terminalen. Skapa filen med koden nedan först.

\begin{minted}{text}
with open("results.csv", "w") as file:
    file.write('Birk,3\n"Rövardotter, Ronja",5\n')
\end{minted}

\begin{itemize}
  \item \textbf{«Filen behöver inte öppnas om modulen \texttt{csv} används»} Skriv ut varje rad från \texttt{csv.\allowbreak{}reader(\allowbreak{}\textquotedbl{}results.\allowbreak{}csv\textquotedbl{})\allowbreak{}} i en \texttt{for}-slinga och läs vad dokumentationen säger om vad \texttt{csv.reader} vill ha.
  \item \textbf{«Läsningen med \texttt{split} avbryts med \texttt{ValueError} på raden med Ronja»} Kör läsningen med \texttt{split}.
  \item \textbf{«Med modulen \texttt{csv} är \texttt{score} ett heltal»} Skriv ut \texttt{type(score)} i slingan med modulen \texttt{csv}.
  \item \textbf{«Med modulen \texttt{csv} får \texttt{name} värdet \texttt{Rövardotter,\allowbreak{} Ronja}, utan citattecken»} Skriv ut \texttt{name} i slingan med modulen \texttt{csv}.
  \item \textbf{«Efter läsningen med modulen \texttt{csv} måste \texttt{file.close()} anropas»} Skriv ut \texttt{file.closed} efter \texttt{with}-blocket.
\end{itemize}

\textbf{I kursen:} vecka 45, föreläsningen \emph{Arbeta med filer}, avsnitten \emph{Öppna och stänga}, \emph{Egna filformat} och \emph{Formatet CSV}.
  \item[Feedback on a right answer] \begin{itemize}
  \item \textbf{«Filen behöver inte öppnas om modulen \texttt{csv} används»} Stämmer inte. \texttt{csv.reader} vill ha något som ger rader, till exempel en öppen fil. En sträng ger ett tecken i taget, så \texttt{csv.\allowbreak{}reader(\allowbreak{}\textquotedbl{}results.\allowbreak{}csv\textquotedbl{})\allowbreak{}} läser tecknen i filnamnet, inte filen.
  \item \textbf{«Läsningen med \texttt{split} avbryts med \texttt{ValueError} på raden med Ronja»} Stämmer. \texttt{split(\textquotedbl{},\textquotedbl{})} vet inget om citattecken, så raden delas vid båda kommatecknen och ger tre delar. Två variabler räcker inte till tre värden.
  \item \textbf{«Med modulen \texttt{csv} är \texttt{score} ett heltal»} Stämmer inte. Alla fält är strängar, så \texttt{score} är \texttt{\textquotedbl{}5\textquotedbl{}}, och dokumentationen säger att ingen typomvandling görs. Du behöver \texttt{int(score)}, precis som med \texttt{split}.
  \item \textbf{«Med modulen \texttt{csv} får \texttt{name} värdet \texttt{Rövardotter,\allowbreak{} Ronja}, utan citattecken»} Stämmer. Citattecknen är en del av filformatet och säger att kommatecknet hör till fältet. Modulen tar bort dem åt dig.
  \item \textbf{«Efter läsningen med modulen \texttt{csv} måste \texttt{file.close()} anropas»} Stämmer inte. \texttt{with} stänger filen när blocket lämnas, oavsett vad som har läst den.
\end{itemize}
  \item[Feedback on a wrong answer] Antingen har du markerat något som inte ska vara med, eller missat något. Du behöver inte gissa här: kör varje prov ovan och jämför vad som händer med vad påståendet säger, ord för ord.
\end{description}

\subsection{Question 13}
\label{ep2:q13}
\emph{Learning objective 5, 10\,000 points, choice.}

\emph{Lärandemål 5: Konstruera interaktiva program}

\textbf{Välj ett alternativ.}

Funktionen \texttt{ask\_grade} ska fråga om och om igen tills användaren har svarat med ett heltal mellan 1 och 10, och sedan returnera talet som ett \texttt{int}. Den får inte krascha, vad användaren än skriver. Vilken version gör det?

The options; the key is marked:
\begin{description}
  \item[{A}] \leavevmode\par
\begin{minted}{python}
def ask_grade():
    grade = int(input("Betyg (1-10): "))
    while not 1 <= grade <= 10:
        print("Betyget ska vara mellan 1 och 10.")
        grade = int(input("Betyg (1-10): "))
    return grade
\end{minted}


    \emph{Feedback:} Inte den här: slingan klarar svar utanför intervallet, men \texttt{int(input(...))} står utan \texttt{try}. Svaret \texttt{sju} avbryter programmet med \texttt{ValueError}, både vid första frågan och i slingan. Att det finns en slinga betyder inte att alla svar är hanterade.
  \item[{B \textit{(key)}}] \leavevmode\par
\begin{minted}{python}
def ask_grade():
    while True:
        try:
            grade = int(input("Betyg (1-10): "))
        except ValueError:
            print("Svara med ett heltal, till exempel 7.")
            continue
        if 1 <= grade <= 10:
            return grade
        print("Betyget ska vara mellan 1 och 10.")
\end{minted}


    \emph{Feedback:} Rätt. Funktionen klarar alla svar, i alla ordningar, och returnerar ett giltigt tal.
  \item[{C}] \leavevmode\par
\begin{minted}{python}
def ask_grade():
    while True:
        try:
            grade = int(input("Betyg (1-10): "))
        except ValueError:
            print("Svara med ett heltal, till exempel 7.")
        if 1 <= grade <= 10:
            return grade
        print("Betyget ska vara mellan 1 och 10.")
\end{minted}


    \emph{Feedback:} Inte den här: svarar du \texttt{sju} först skrivs meddelandet ut, och sedan går funktionen vidare till \texttt{if}, där \texttt{grade} inte har något värde: \texttt{UnboundLocalError}. Svarar du \texttt{11} och sedan \texttt{sju} kraschar den inte, men du får två meddelanden. Efter ett \texttt{except} fortsätter programmet på nästa rad.
  \item[{D}] \leavevmode\par
\begin{minted}{python}
def ask_grade():
    try:
        grade = int(input("Betyg (1-10): "))
        if 1 <= grade <= 10:
            return grade
        print("Betyget ska vara mellan 1 och 10.")
    except ValueError:
        print("Svara med ett heltal, till exempel 7.")
    ask_grade()
\end{minted}


    \emph{Feedback:} Inte den här: funktionen anropar sig själv för att fråga igen, men svaret kommer aldrig tillbaka till den första anroparen. Svarar du \texttt{7} direkt blir det rätt, men \texttt{sju} och sedan \texttt{7} ger \texttt{None}. Inget kraschar, så du måste titta på returvärdet.
\end{description}

The item's feedback:
\begin{description}
  \item[General feedback] \textbf{Tänk på:} användaren kan svara på två felaktiga sätt, med något som inte är ett heltal och med ett heltal utanför intervallet. Funktionen ska klara båda, i vilken ordning som helst, och till sist lämna tillbaka talet.

\textbf{Prova:} lägg in en version i ett program, skriv \texttt{print(\allowbreak{}ask\_\allowbreak{}grade(\allowbreak{})\allowbreak{})\allowbreak{}} sist och kör. Svara \texttt{sju}, sedan \texttt{11} och sist \texttt{7} i dialogrutan. Kör en gång till och svara \texttt{7} direkt. Sista utskriften är returvärdet, och den ska båda gångerna vara \texttt{7}.

\textbf{I kursen:} vecka 40, föreläsningen \emph{Upprepningar}, avsnitten \emph{Inmatning som frågar om igen} och \emph{Allt på en gång: ett gissningsspel}; vecka 39, föreläsningen \emph{Villkor och styrstrukturer}, avsnittet \emph{Att fråga om igen}; Laboration 3.
  \item[Feedback on a right answer] Slingan står för omfrågningen, \texttt{try} för svar som inte går att omvandla, och \texttt{continue} hoppar över kontrollen när det inte finns något tal att kontrollera. \texttt{sju} fångas och nästa varv börjar direkt. \texttt{11} passerar omvandlingen, stoppas av kontrollen och får ett meddelande. \texttt{7} returneras som \texttt{int}. Inget svar kan avbryta funktionen.
  \item[Feedback on a wrong answer] Felen syns bara för vissa svar i viss ordning. Pröva din version med text som allra första svar, och med text efter ett tal utanför intervallet. Titta på returvärdet, inte bara på det som händer i dialogrutan.
\end{description}

\subsection{Question 14}
\label{ep2:q14}
\emph{Learning objective 5, 10\,000 points, choice.}

\emph{Lärandemål 5: Konstruera interaktiva program}

\textbf{Välj ett alternativ.}

\begin{minted}{python}
"""Räknar ut ett pris med moms."""

VAT_RATE = 0.25


def with_vat(price):
    """Returnerar price med moms."""
    return price * (1 + VAT_RAET)


def main():
    """Frågar efter ett pris och skriver ut det med moms."""
    try:
        price = float(input("Pris utan moms (kr): "))
        print(f"Med moms: {with_vat(price):.2f} kr")
    except Exception:
        print("Priset måste skrivas med siffror, till exempel 99.50.")


if __name__ == "__main__":
    main()
\end{minted}

Användaren svarar \texttt{100} och får ändå meddelandet att priset måste skrivas med siffror. Vad är orsaken, och vad bör programmeraren ändra för att sådana fel ska synas i fortsättningen?

The options; the key is marked:
\begin{description}
  \item[{A}] \texttt{input} lämnar tillbaka en sträng, så \texttt{price *\allowbreak{} (\allowbreak{}1 +\allowbreak{} VAT\_\allowbreak{}RAET)\allowbreak{}} ger \texttt{TypeError}. Omvandlingen måste läggas inne i \texttt{with\_vat}.

    \emph{Feedback:} Inte den här: det är ett vanligt fel, men inte här. \texttt{float(...)} står redan runt \texttt{input}, så \texttt{price} är ett flyttal när \texttt{with\_vat} anropas, och det blir inget \texttt{TypeError}.
  \item[{B}] \texttt{float} kan inte omvandla \texttt{100}, eftersom talet saknar decimalpunkt. Byt till \texttt{int} för att också ta emot heltal.

    \emph{Feedback:} Inte den här: \texttt{float(\textquotedbl{}100\textquotedbl{})} ger \texttt{100.0}, så ett tal utan decimalpunkt går bra. Prova det i terminalen. Ett meddelande är en text någon skrev, och det kan peka fel.
  \item[{C \textit{(key)}}] \texttt{VAT\_RAET} är felstavat, och \texttt{except Exception} fångar också det \texttt{NameError} det ger. Rätta stavningen och fånga bara \texttt{ValueError}.

    \emph{Feedback:} Rätt. Både orsaken och ändringen stämmer.
  \item[{D}] \texttt{:.2f} går inte att använda på ett funktionsanrop i en f-sträng, och det ger \texttt{ValueError}. Lägg resultatet i en variabel först.

    \emph{Feedback:} Inte den här: \texttt{:.2f} fungerar på vilket uttryck som helst, också ett funktionsanrop. Kursen skriver till exempel \texttt{\{bmi(\allowbreak{}weight,\allowbreak{} length)\allowbreak{}:.\allowbreak{}1f\}}. Det behövs ingen extra variabel.
\end{description}

The item's feedback:
\begin{description}
  \item[General feedback] \textbf{Tänk på:} meddelandet säger vad programmeraren trodde skulle gå fel. Det säger inte vad som faktiskt gick fel. Frågan har två delar, orsaken och vad som ska ändras.

\textbf{Prova:} kör programmet utan \texttt{try} och \texttt{except}, med enbart de två raderna inuti, och svara \texttt{100}. Läs sedan hela felmeddelandet. Pröva också den ändring ditt alternativ föreslår och se om problemet försvinner.

\textbf{I kursen:} vecka 39, föreläsningen \emph{Inmatning och felhantering}, avsnitten \emph{Flera sorters fel} och \emph{Ett begripligt meddelande}.
  \item[Feedback on a right answer] \texttt{VAT\_RAET} är felstavat, och \texttt{except Exception} fångar också det \texttt{NameError} det ger. Användaren får då skulden för ett fel i programmet. Fånga bara det du kan svara på, här \texttt{ValueError} från \texttt{float}. Rättar du stavningen och fångar bara \texttt{ValueError} får användaren rätt besked, och nästa programmeringsfel syns direkt.
  \item[Feedback on a wrong answer] Jämför förklaringen du valde med felmeddelandet du får utan \texttt{try} och \texttt{except}. Stämmer sorten av fel, och stämmer raden?
\end{description}

\subsection{Question 15}
\label{ep2:q15}
\emph{Learning objective 5, 10\,000 points, choice.}

\emph{Lärandemål 5: Konstruera interaktiva program}

\textbf{Välj ett alternativ.}

\begin{minted}{python}
"""Säljer biljetter."""

PRICE = 100


def main():
    """Frågar efter antalet biljetter och skriver ut priset."""
    count = input("Hur många biljetter? ")
    if count == 2:
        print("Två biljetter: du får en popcorn på köpet.")
    print(f"Att betala: {count * PRICE} kr")


if __name__ == "__main__":
    main()
\end{minted}

När användaren svarar \texttt{2} ska programmet skriva ut raden om popcorn och \texttt{Att betala: 200 kr}. Det gör det inte. Vilken enda ändring räcker för att båda raderna ska bli rätt?

The options; the key is marked:
\begin{description}
  \item[{A \textit{(key)}}] Skriv \texttt{count =\allowbreak{} int(\allowbreak{}input(\allowbreak{}\textquotedbl{}Hur många biljetter?\allowbreak{} \textquotedbl{})\allowbreak{})\allowbreak{}} på raden som frågar.

    \emph{Feedback:} Rätt. En ändring på rätt ställe rättar båda raderna.
  \item[{B}] Skriv \texttt{if count == \textquotedbl{}2\textquotedbl{}:} i stället för \texttt{if count == 2:}.

    \emph{Feedback:} Inte den här: popcornraden kommer, men priset blir hundra tvåor i rad, inte 200. Du har åtgärdat ett symptom, och det andra är kvar.
  \item[{C}] Skriv \texttt{\{int(\allowbreak{}count *\allowbreak{} PRICE)\allowbreak{}\}} i stället för \texttt{\{count * PRICE\}} i f-strängen.

    \emph{Feedback:} Inte den här: \texttt{count * PRICE} är redan hundra tvåor i rad när \texttt{int} får dem, så det blir ett tal med hundra siffror. Popcornraden saknas dessutom fortfarande.
  \item[{D}] Skriv \texttt{PRICE = \textquotedbl{}100\textquotedbl{}} i stället för \texttt{PRICE = 100}.

    \emph{Feedback:} Inte den här: två strängar går inte att multiplicera, så programmet avbryts med \texttt{TypeError}, och popcornraden saknas fortfarande. Att göra typerna lika är inte målet.
\end{description}

The item's feedback:
\begin{description}
  \item[General feedback] \textbf{Tänk på:} två saker blir fel. Har de en gemensam orsak? Frågan söker en enda ändring som rättar båda.

\textbf{Prova:} skriv ut \texttt{type(count)} direkt efter raden med \texttt{input}, och kör \texttt{print(\textquotedbl{}2\textquotedbl{} == 2)} och \texttt{print(\textquotedbl{}2\textquotedbl{} * 3)}. Pröva sedan varje ändring för sig: kör programmet och svara \texttt{2} i dialogrutan.

\textbf{I kursen:} vecka 39, föreläsningen \emph{Inmatning och felhantering}, avsnitten \emph{Inmatning} och \emph{Datatyper och operationer}.
  \item[Feedback on a right answer] \texttt{input} lämnar alltid tillbaka en sträng. Därför är \texttt{\textquotedbl{}2\textquotedbl{} == 2} falskt, och \texttt{\textquotedbl{}2\textquotedbl{} * 100} blir hundra tvåor i rad. Två symptom har en orsak. Med \texttt{int} runt \texttt{input} är \texttt{count} talet 2, så \texttt{count == 2} är sant och \texttt{count * PRICE} blir 200. Omvandla vid inläsningen, så innehåller variabeln det den heter. Läs mer i föreläsningen \emph{Inmatning och felhantering}, avsnitten \emph{Vad är det vi får tillbaka?} och \emph{Var hör omvandlingen hemma?}.
  \item[Feedback on a wrong answer] Kör programmet med ändringen du valde och svara \texttt{2}. Vilken av de två raderna är fortfarande fel, eller vad händer i stället? Fråga dig vad \texttt{count} innehåller när den raden körs.
\end{description}

\subsection{Question 16}
\label{ep2:q16}
\emph{Learning objective 6, 100\,000 points, choice.}

\emph{Lärandemål 6: Använda och konstruera sammansatta datatyper och klasser}

\textbf{Välj ett alternativ.}

Skriv ett litet program i terminalen som löser uppgiften, och välj sedan resultatet.

Siffersumman av ett heltal är summan av dess siffror: 123 har siffersumman 6. Gå igenom alla heltal från och med 1 till och med 500, och räkna med en uppslagslista hur många av talen som har varje siffersumma. Vilken siffersumma är vanligast, och hur många av talen har den?

The options; the key is marked:
\begin{description}
  \item[{A}] Siffersumman 1, som 50 tal har.

    \emph{Feedback:} Inte den här: det här får du med \texttt{total =\allowbreak{} int(\allowbreak{}digit)\allowbreak{}} i stället för \texttt{total +\allowbreak{}=\allowbreak{} int(\allowbreak{}digit)\allowbreak{}}. Varje varv skriver då över summan, så bara sista siffran är kvar, och 123 ger 3 i stället för 6. Alla tio siffror får 50 tal, och 1 kommer först.
  \item[{B}] Siffersumman 9, som 49 tal har.

    \emph{Feedback:} Inte den här: det här får du om siffersumman räknas i en slinga med \texttt{while number \textgreater{} 9:} och \texttt{\%} och \texttt{//}. Den första siffran kommer då aldrig med, så 123 ger 5 i stället för 6.
  \item[{C}] Siffersumman 11, som 43 tal har.

    \emph{Feedback:} Inte den här: kontrollera räknaren. Gå igenom bara talen 1 till 20, skriv ut hela uppslagslistan och jämför med vad du räknar för hand: hur många av talen har siffersumman 2?
  \item[{D \textit{(key)}}] Siffersumman 11, som 44 tal har.

    \emph{Feedback:} Rätt. Siffersumman, räknaren och sökningen efter det största värdet stämmer alla.
  \item[{E}] Siffersumman 11, som 45 tal har.

    \emph{Feedback:} Inte den här: kontrollera räknaren. Gå igenom bara talen 1 till 20, skriv ut hela uppslagslistan och jämför med vad du räknar för hand: hur många av talen har siffersumman 2?
  \item[{F}] Siffersumman 22, som 1 tal har.

    \emph{Feedback:} Inte den här: 22 är vad \texttt{max(counts)} ger. \texttt{max} och \texttt{for} över en uppslagslista arbetar med nycklarna, så du får den största siffersumman (499 ger 22), inte den vanligaste.
\end{description}

The item's feedback:
\begin{description}
  \item[General feedback] \textbf{Tänk på:} programmet har tre delar: att räkna ut en siffersumma, att räkna hur många tal som har varje siffersumma, och att hitta den vanligaste. Testa dem var för sig.

\textbf{Prova:} kontrollera först siffersumman: 123 ska ge 6, 500 ska ge 5 och 9 ska ge 9. Gå sedan igenom bara talen 1 till 20, skriv ut hela uppslagslistan och jämför med vad du räknar för hand. Skriv till sist ut vilken nyckel och vilket värde du väljer, och kontrollera att det värdet verkligen är störst. Tar programmet aldrig slut: tryck på Starta om miljön eller ladda om sidan. Dina filer försvinner då.

\textbf{I kursen:} vecka 41, föreläsningen \emph{Behållare: Uppslagslistor}, avsnitten \emph{Uppslagslistor}, \emph{När nyckeln saknas} och \emph{Att gå igenom en uppslagslista}; vecka 40, föreläsningen \emph{Upprepningar}, avsnittet \emph{Siffersumman, tre gånger}; vecka 44, \emph{Övning: Fler behållare och mer om klasser}, uppgiften \emph{Bokstavsräknaren}.
  \item[Feedback on a right answer] Uppslagslistan har siffersumman som nyckel och antalet tal som värde, och \texttt{get} med reservvärdet 0 sköter första gången en siffersumma dyker upp. \texttt{key=counts.get} gör att \texttt{max} jämför antalen, inte nycklarna. Siffersumman 11 har 44 tal, och närmast kommer 10 och 12 med 43 tal var. Så här kan det se ut:

\begin{minted}{python}
def digit_sum(number):
    """Returnerar siffersumman av number."""
    total = 0
    for digit in str(number):
        total += int(digit)
    return total


def main():
    """Skriver ut den vanligaste siffersumman bland talen 1 till 500."""
    counts = {}
    for number in range(1, 501):
        total = digit_sum(number)
        counts[total] = counts.get(total, 0) + 1
    most_common = max(counts, key=counts.get)
    print(most_common, counts[most_common])


if __name__ == "__main__":
    main()
\end{minted}
  \item[Feedback on a wrong answer] Vilken av de tre delarna brister? Skriv ut siffersumman för några tal du kan räkna för hand, skriv ut uppslagslistan, och fråga dig om du jämför nycklar eller värden när du letar efter det vanligaste.
\end{description}

\subsection{Question 17}
\label{ep2:q17}
\emph{Learning objective 6, 100\,000 points, choice.}

\emph{Lärandemål 6: Använda och konstruera sammansatta datatyper och klasser}

\textbf{Välj ett alternativ.}

\begin{minted}{python}
class Student:
    """En student med namn och personnummer."""

    def __init__(self, name, pnr):
        self.__name = name
        self.__pnr = pnr

    def get_pnr(self):
        """Returnerar personnumret."""
        return self.__pnr

    def __str__(self):
        return f"Namn: {self.__name} Personnr: {self.__pnr}"


class School:
    """En skola med sina studenter."""

    def __init__(self):
        self.__students = {}

    ...  # Här ska metoderna add och find stå.
\end{minted}

Metoden \texttt{add} ska spara ett \texttt{Student}-objekt i skolan. Metoden \texttt{find} ska ta ett personnummer och returnera det \texttt{Student}-objektet, eller \texttt{None} om ingen student har det numret. Vilka metoder gör det?

The options; the key is marked:
\begin{description}
  \item[{A}] \leavevmode\par
\begin{minted}{python}
def add(self, student):
    self.__students[student.__pnr] = student

def find(self, pnr):
    return self.__students.get(pnr)
\end{minted}


    \emph{Feedback:} Inte den här: \texttt{\_\_pnr} är privat för \texttt{Student}. Inuti \texttt{School} skriver Python om \texttt{student.\_\_pnr} till \texttt{student.\allowbreak{}\_\allowbreak{}School\_\allowbreak{}\_\allowbreak{}pnr}, och det attributet finns inte, så \texttt{add} ger \texttt{AttributeError}. Det är just det inkapsling innebär.
  \item[{B}] \leavevmode\par
\begin{minted}{python}
def add(self, student):
    students[student.get_pnr()] = student

def find(self, pnr):
    return students.get(pnr)
\end{minted}


    \emph{Feedback:} Inte den här: utan \texttt{self.} är \texttt{students} en vanlig variabel, och någon sådan finns inte i metoden, så det blir \texttt{NameError}. Ett attribut är inte en variabel som gäller i hela klassen.
  \item[{C}] \leavevmode\par
\begin{minted}{python}
def add(self, student):
    self.__students[student] = student.get_pnr()

def find(self, pnr):
    return self.__students.get(pnr)
\end{minted}


    \emph{Feedback:} Inte den här: nyckel och värde har bytt plats, så \texttt{get} ger alltid \texttt{None} när du söker på numret, och inget kraschar. Tänk från uppslaget: vad har du när du frågar?
  \item[{D \textit{(key)}}] \leavevmode\par
\begin{minted}{python}
def add(self, student):
    self.__students[student.get_pnr()] = student

def find(self, pnr):
    return self.__students.get(pnr)
\end{minted}


    \emph{Feedback:} Rätt. Båda metoderna gör det som efterfrågas, också när numret saknas.
\end{description}

The item's feedback:
\begin{description}
  \item[General feedback] \textbf{Tänk på:} tre saker ska stämma. Hur når en metod objektets eget attribut? Hur når \texttt{School} något som är privat i \texttt{Student}? Vad har du när du anropar \texttt{find}, och vad vill du få tillbaka?

\textbf{Prova:} skriv klasserna med en version av metoderna och lägg koden nedan sist i programmet. Den ska först skriva ut studenten och sedan \texttt{None}.

\begin{minted}{python}
school = School()
school.add(Student("Ronja", "1234"))
print(school.find("1234"))
print(school.find("9999"))
\end{minted}

\textbf{I kursen:} vecka 41, föreläsningen \emph{Klasser och objekt}, avsnittet \emph{Metoder och attribut}; föreläsningen \emph{Behållare: Uppslagslistor}, avsnittet \emph{När nyckeln saknas}; vecka 44, \emph{Övning: Fler behållare och mer om klasser}, uppgiften \emph{Banken}; Laboration 5.
  \item[Feedback on a right answer] Tre saker stämmer: attributet nås genom \texttt{self}, ett annat objekts privata attribut nås genom dess läsmetod \texttt{get\_pnr}, och nyckeln är personnumret, eftersom det är vad du har när du frågar. \texttt{find} slår upp numret och får tillbaka objektet, och \texttt{get} returnerar \texttt{None} när nyckeln saknas.
  \item[Feedback on a wrong answer] Kör din version med koden ovan. Avbryts programmet, så läs vilket namn felmeddelandet nämner. Får du \texttt{None} båda gångerna, så skriv ut uppslagslistan och titta på vad som är nyckel och vad som är värde.
\end{description}

\subsection{Question 18}
\label{ep2:q18}
\emph{Learning objective 6, 100\,000 points, choice.}

\emph{Lärandemål 6: Använda och konstruera sammansatta datatyper och klasser}

\textbf{Välj ett alternativ.}

Ett brädspel ska hålla reda på vilken pjäs som står på vilken ruta. En ruta anges med rad och kolumn. Programmet ska många gånger fråga vilken pjäs som står på en viss ruta, utan att gå igenom alla pjäser, och pjäserna flyttas under spelets gång. Vilket sätt att lagra pjäserna passar bäst?

The options; the key is marked:
\begin{description}
  \item[{A}] \leavevmode\par
\begin{minted}{python}
board = {[0, 4]: "kung", [7, 4]: "drottning"}
\end{minted}


    \emph{Feedback:} Inte den här: en lista kan ändras och får därför inte vara nyckel. Redan raden som skapar \texttt{board} ger \texttt{TypeError}.
  \item[{B \textit{(key)}}] \leavevmode\par
\begin{minted}{python}
board = {(0, 4): "kung", (7, 4): "drottning"}
\end{minted}


    \emph{Feedback:} Rätt. Formen klarar både uppslaget och flytten.
  \item[{C}] \leavevmode\par
\begin{minted}{python}
board = [(0, 4, "kung"), (7, 4, "drottning")]
\end{minted}


    \emph{Feedback:} Inte den här: det går att lagra pjäserna så, men för att hitta pjäsen på en ruta måste du gå igenom listan, och det är det uppgiften säger att du ska slippa. Lista är ofta förstahandsvalet, men här passar den inte operationerna.
  \item[{D}] \leavevmode\par
\begin{minted}{python}
board = {"kung": (0, 4), "drottning": (7, 4)}
\end{minted}


    \emph{Feedback:} Inte den här: den svarar på var en viss pjäs står, inte på vilken pjäs som står på en ruta. Två pjäser av samma slag går inte heller att lagra, eftersom samma nyckel skrivs över.
\end{description}

The item's feedback:
\begin{description}
  \item[General feedback] \textbf{Tänk på:} välj behållare efter vilka operationer programmet behöver. Vilken fråga ska programmet ställa många gånger, och vad vet det när det frågar?

\textbf{Prova:} skriv för varje form två rader kod: en som tar reda på vilken pjäs som står på rutan (0, 4), och en som flyttar den till rutan (1, 4). Kör dem. Går koden att köra alls? Hur många pjäser måste den gå igenom?

\textbf{I kursen:} vecka 41, föreläsningen \emph{Behållare: Uppslagslistor}, avsnitten \emph{Vad duger som nyckel?} och \emph{Att välja behållare}; vecka 40, föreläsningen \emph{Behållare: Tupler}, avsnittet \emph{Tupler}.
  \item[Feedback on a right answer] Programmet frågar «vad står på den här rutan?», så rutan ska vara nyckeln. En tupel duger som nyckel eftersom den inte kan ändras, och uppslaget går direkt, utan att du går igenom andra pjäser. Pjäserna kan ändå flyttas, eftersom du kan lägga till och ta bort nycklar:

\begin{minted}{python}
board = {(0, 4): "kung", (7, 4): "drottning"}
print(board.get((0, 4)))
board[(1, 4)] = "kung"
del board[(0, 4)]
print(board)
\end{minted}
  \item[Feedback on a wrong answer] Skriv raden som svarar på vilken pjäs som står på rutan (0, 4) med den form du valde. Går den att köra? Klarar den frågan utan att gå igenom alla pjäser?
\end{description}

\subsection{Question 19}
\label{ep2:q19}
\emph{Learning objective 7, 1\,000\,000 points, choice.}

\emph{Lärandemål 7: Dela upp ett program}

\textbf{Välj ett alternativ.}

\begin{minted}{python}
def report_total(prices):
    """Skriver ut summan av prices med moms."""
    total = 0
    for price in prices:
        total += price * 1.25
    print(f"Summa med moms: {total:.2f} kr")


def report_student_total(prices):
    """Skriver ut summan av prices med moms och studentrabatt."""
    total = 0
    for price in prices:
        total += price * 1.25
    print(f"Studentpris med moms: {total * 0.9:.2f} kr")
\end{minted}

Programmet ska byggas ut. Summorna ska också kunna skrivas till en fil, utan någon utskrift i terminalen, och momssatsen ändras ibland. Vilken uppdelning passar bäst, med tanke på principerna DRY och SRP?

The options; the key is marked:
\begin{description}
  \item[{A}] Slå ihop de två till en funktion \texttt{report\_\allowbreak{}total(\allowbreak{}prices,\allowbreak{} student)\allowbreak{}}, som med en \texttt{if}-sats väljer rabatt och sedan skriver ut summan. Då finns bara en funktion att ändra när momsen eller utskriften ändras.

    \emph{Feedback:} Inte den här: upprepningen försvinner, och det är bra enligt DRY, men räkning och utskrift sitter ihop i samma funktion. Den som vill skriva summan till en fil får den också utskriven i terminalen. Funktionen har två ansvar.
  \item[{B}] Bryt ut utskriften till en funktion \texttt{print\_\allowbreak{}total(\allowbreak{}text,\allowbreak{} total)\allowbreak{}}, som båda funktionerna anropar. Då står utskriften på ett ställe, medan slingorna och talet 1.25 står kvar där de är.

    \emph{Feedback:} Inte den här: det ser ut som SRP, men det löser inte ändringarna. Slingorna och talet 1.25 står kvar på två ställen, så en ny momssats kräver två ändringar, och uträkningen sitter kvar i funktioner som skriver ut.
  \item[{C}] Flytta de två funktionerna oförändrade till en egen modul \texttt{reports.py} och importera den där de behövs. Då kan alla program som behöver summorna använda samma kod i stället för att ha egna kopior.

    \emph{Feedback:} Inte den här: moduler är ett bra sätt att återanvända kod, men funktionerna ändras inte av att flyttas. Upprepningen följer med, och de skriver fortfarande ut.
  \item[{D \textit{(key)}}] Inför konstanten \texttt{VAT\_RATE} och en funktion \texttt{total\_\allowbreak{}with\_\allowbreak{}vat(\allowbreak{}prices)\allowbreak{}} som returnerar summan med moms utan att skriva ut något. Den som anropar bestämmer vad som görs med summan: drar av rabatten, skriver ut den eller sparar den i en fil.

    \emph{Feedback:} Rätt. Förslaget klarar båda ändringarna utan att någon funktion behöver skrivas om.
\end{description}

The item's feedback:
\begin{description}
  \item[General feedback] \textbf{Tänk på:} frågan nämner två ändringar. Momssatsen ändras ibland, och summorna ska också kunna skrivas till en fil utan någon utskrift i terminalen. Genomför båda i tanken på varje förslag. Hur många ställen måste du ändra (DRY)? Hur många uppgifter har varje funktion (SRP)?

\textbf{Prova:} gör samma sak med ett eget program: vad måste du skriva om när en del av det ändras?

\textbf{I kursen:} vecka 38, föreläsningen \emph{Funktioner}, avsnitten \emph{Att dela upp ett program i funktioner} och \emph{Tre principer med namn} (vad DRY och SRP vinner när programmet ändras); \emph{Övning: Funktioner och variabler}, uppgiften \emph{Kakelfabriken}. Bedömning: rättningsmatrisen i \emph{Laboration (1) funktioner, variabler och utskrifter}, kriterierna \emph{Uppdelning i funktioner/metoder}, \emph{Ingen kodupprepning} och \emph{Parametrar/returvärden}.
  \item[Feedback on a right answer] Momssatsen står på ett ställe, och uträkningen på ett ställe (DRY). Funktionen gör en sak, räknar, och returnerar svaret (SRP). Den som anropar bestämmer sedan om summan ska få rabatt, skrivas ut eller sparas i en fil.

\begin{minted}{python}
VAT_RATE = 0.25


def total_with_vat(prices):
    """Returnerar summan av prices med moms."""
    total = 0
    for price in prices:
        total += price * (1 + VAT_RATE)
    return total
\end{minted}

\texttt{print(\allowbreak{}f\textquotedbl{}Studentpris med moms: \{total\_\allowbreak{}with\_\allowbreak{}vat(\allowbreak{}prices)\allowbreak{} *\allowbreak{} 0.\allowbreak{}9:.\allowbreak{}2f\} kr\textquotedbl{})\allowbreak{}}
  \item[Feedback on a wrong answer] Ta förslaget du valde och skriv upp vad som måste göras för varje ändring: en ny momssats, och en summa i en fil utan utskrift. Går båda utan att du ändrar på flera ställen eller skriver om funktionerna?
\end{description}

\subsection{Question 20}
\label{ep2:q20}
\emph{Learning objective 7, 1\,000\,000 points, choice.}

\emph{Lärandemål 7: Dela upp ett program}

\textbf{Välj ett alternativ.}

\begin{minted}{python}
"""Räknar ut moms på ett pris."""


def vat(price):
    """Returnerar momsen på price."""
    return price * rate


def main():
    """Skriver ut momsen på 100 kronor."""
    rate = 0.25
    print(vat(100))


if __name__ == "__main__":
    main()
\end{minted}

Vad händer när programmet körs, och varför?

The options; the key is marked:
\begin{description}
  \item[{A}] Programmet skriver ut \texttt{25.0}, eftersom \texttt{main} har skapat \texttt{rate} innan \texttt{vat} anropas. Ingen ändring behövs.

    \emph{Feedback:} Inte den här: \texttt{main} har visserligen skapat \texttt{rate} när \texttt{vat} körs, men det räcker inte för att \texttt{vat} ska kunna använda det. Kör programmet så ser du vad som händer. Att ett namn finns just nu betyder inte att det går att nå överallt.
  \item[{B}] Programmet avbryts med \texttt{NameError}, eftersom \texttt{def vat} står före raden \texttt{rate = 0.25} i filen. Det löses genom att flytta \texttt{def vat} till efter \texttt{def main}.

    \emph{Feedback:} Inte den här: ordningen mellan funktionerna spelar ingen roll här, eftersom båda redan finns när \texttt{main()} körs. Flytta \texttt{def vat} och kör: ingenting ändras. Det är en annan regel du tänker på, den om att en funktion måste vara definierad innan den anropas.
  \item[{C \textit{(key)}}] Programmet avbryts med \texttt{NameError}, eftersom \texttt{rate} är lokal i \texttt{main} och en funktion aldrig letar bland namnen i den funktion som anropade den. Det löses genom att ge \texttt{vat} en parameter \texttt{rate} och skicka med värdet.

    \emph{Feedback:} Rätt. Både utfallet, orsaken och lösningen stämmer.
  \item[{D}] Programmet avbryts med \texttt{TypeError}, eftersom \texttt{vat} anropas med ett argument men använder två värden. Det löses genom att anropa \texttt{vat(100, 0.25)} utan att ändra funktionshuvudet.

    \emph{Feedback:} Inte den här: \texttt{vat(100)} stämmer med \texttt{def vat(price)}, så anropet går bra. Python jämför inte kroppen med funktionshuvudet. Med \texttt{vat(100, 0.25)} utan nytt funktionshuvud får du däremot \texttt{TypeError}, eftersom funktionen bara tar ett argument.
\end{description}

The item's feedback:
\begin{description}
  \item[General feedback] \textbf{Tänk på:} varje alternativ säger tre saker: vad som händer, varför, och vad som löser det. Alla tre ska stämma.

\textbf{Prova:} kör programmet och läs felmeddelandet, om det kommer något. Gör sedan den ändring ett alternativ föreslår och kör igen. Skriv också upp vilka namn varje funktion kan se.

\textbf{I kursen:} vecka 38, föreläsningen \emph{Funktioner}, avsnittet \emph{Parametrar, argument och namn} och där \emph{Huvudprogrammet i en egen funktion} (varför en funktion ska få sina värden som parametrar). Bedömning: rättningsmatrisen i \emph{Laboration (1) funktioner, variabler och utskrifter}, kriteriet \emph{Parametrar/returvärden}.
  \item[Feedback on a right answer] \texttt{rate} finns bara inne i \texttt{main}. En funktion ser sina egna namn och namnen på filens översta nivå, men aldrig namnen i den funktion som anropade den, så \texttt{vat} hittar inget med det namnet och programmet avbryts med \texttt{NameError}. Det en funktion behöver ska den få som parameter:

\begin{minted}{python}
def vat(price, rate):
    """Returnerar momsen på price med momssatsen rate."""
    return price * rate


def main():
    """Skriver ut momsen på 100 kronor."""
    rate = 0.25
    print(vat(100, rate))


if __name__ == "__main__":
    main()
\end{minted}

Läs mer i föreläsningen \emph{Funktioner}, avsnitten \emph{Funktionens namn är dess egna} och \emph{Huvudprogrammet i en egen funktion}.
  \item[Feedback on a wrong answer] Pröva alternativet du valde: stämmer det som händer när du kör programmet, och löser ändringen problemet? Båda delarna behöver hålla.
\end{description}

\subsection{Question 21}
\label{ep2:q21}
\emph{Learning objective 7, 1\,000\,000 points, choice.}

\emph{Lärandemål 7: Dela upp ett program}

\textbf{Välj ett alternativ.}

Ett program är uppdelat på två filer i samma katalog. Filen \texttt{prices.py}:

\begin{minted}{python}
"""Prisberäkningar med moms."""

print(f"Modulen heter {__name__}")


def with_vat(price):
    """Returnerar price med 25 procents moms."""
    return price * 1.25


def main():
    """Provkör with_vat."""
    print(with_vat(100))


if __name__ == "__main__":
    main()
\end{minted}

och filen \texttt{shop.py}:

\begin{minted}{python}
"""Räknar ut priset i en butik."""

import prices


def main():
    """Skriver ut priset för en vara."""
    print(prices.with_vat(200))


if __name__ == "__main__":
    main()
\end{minted}

Vad skrivs ut när man kör \texttt{python3 shop.py}?

The options; the key is marked:
\begin{description}
  \item[{A \textit{(key)}}] \leavevmode\par
\begin{minted}{text}
Modulen heter prices
250.0
\end{minted}


    \emph{Feedback:} Rätt. Importen, villkoret och anropet i \texttt{shop.py} ger exakt de här raderna.
  \item[{B}] \leavevmode\par
\begin{minted}{text}
250.0
\end{minted}


    \emph{Feedback:} Inte den här: \texttt{import} gör mer än att göra namnen tillgängliga, den kör filen. Titta på vilka rader som står på översta nivån i \texttt{prices.py}.
  \item[{C}] \leavevmode\par
\begin{minted}{text}
Modulen heter prices
125.0
250.0
\end{minted}


    \emph{Feedback:} Inte den här: \texttt{main()} i \texttt{prices.py} står bakom ett villkor. Att filen körs vid en import betyder inte att allt i den körs. Vilket värde har \texttt{\_\_name\_\_} när filen importeras?
  \item[{D}] \leavevmode\par
\begin{minted}{text}
Modulen heter __main__
125.0
250.0
\end{minted}


    \emph{Feedback:} Inte den här: det här liknar vad du skulle få om du startade \texttt{prices.py} direkt, men här är det \texttt{shop.py} som startas. Värdet på \texttt{\_\_name\_\_} beror på hur filen körs.
\end{description}

The item's feedback:
\begin{description}
  \item[General feedback] \textbf{Tänk på:} tre saker avgör svaret. Vad gör \texttt{import} med filen? Vilket värde har \texttt{\_\_name\_\_} i en fil som importeras? Vilka \texttt{main} anropas?

\textbf{Prova:} kommandot \texttt{python3 shop.py} går inte att köra i webbläsarens terminal, men du kan se vad en import gör genom att köra koden nedan och läsa vad som skrivs ut. Filen körs bara vid första importen, och ändrar du modulen och importerar den igen får du den gamla versionen. Tryck då på Starta om miljön, men tänk på att dina filer också försvinner.

\begin{minted}{python}
with open("mymodule.py", "w") as file:
    file.write('print(f"Modulen heter {__name__}")\n')
    file.write('if __name__ == "__main__":\n')
    file.write('    print("Villkoret är sant")\n')

import mymodule
\end{minted}

\textbf{I kursen:} vecka 40, föreläsningen \emph{Moduler och paket}, avsnitten \emph{En egen modul} och \emph{Vad händer när en modul importeras?}.
  \item[Feedback on a right answer] \texttt{import prices} kör filen uppifrån och ner. Då körs \texttt{print} på översta nivån, och \texttt{\_\_name\_\_} är \texttt{prices}, så \texttt{Modulen heter prices} skrivs ut. Villkoret sist i filen är falskt, så provkörningen körs inte. Sedan skriver \texttt{main} i \texttt{shop.py} ut \texttt{prices.\allowbreak{}with\_\allowbreak{}vat(\allowbreak{}200)\allowbreak{}}, alltså \texttt{250.0}.
  \item[Feedback on a wrong answer] Gå igenom \texttt{prices.py} rad för rad som om den just importerats. Vilka rader körs, och vilket värde har \texttt{\_\_name\_\_}? Gör sedan samma sak med \texttt{shop.py}.
\end{description}
\end{otherlanguage}
